\section{A coefficient-free Boolean compiler}
\begin{lemma}\label{lem:raw-Boolean-tanh}
There is a numerical $b\ge2$ such that the following holds.  Let $g$ be any
ReLU network of width $C\ge1$ and hidden depth $h\ge1$, with arbitrary
finite real coefficients and affine biases.  On $\{-1,1\}^d$, for every
$0<\epsilon<1$ and every $H\ge b(h+1)$ there is a tanh network $q$ of
hidden depth $H$, with a raw affine readout, satisfying
\[
 \sup_x|q(x)-\operatorname{clip}(g(x),[-1,1])|\le\epsilon,
 \qquad |q(x)|\le1+\epsilon,
\]
and
\[
 m,B\le\left[4(d+1)(C+1)^h
              +\lceil\log_2(32/\epsilon)\rceil+2\right]^b.
\]
The coefficient and width bounds do not grow with the additional learner
depth $H$.  On the entire input cube the same raw implementation satisfies
$|q|\le4$, because its signed-bit readout has coefficient sum below two
and constant offset of magnitude below two.
\end{lemma}
\begin{proof}
Use the coefficient-free unweighted-majority circuit compiler of
\cite{booleancompiler}. Its precise input is an arbitrary-real ReLU
network on $d$ Boolean coordinates, of width $C$ and depth $h$; it outputs
a numerical-constant-depth-per-layer threshold circuit whose size is
polynomial in $4(d+1)(C+1)^h+\log(1/\epsilon)$, independently of all
original coefficient magnitudes.  The construction lifts input-to-node
path products to independent coefficients, fixes the finite Boolean
activation patterns, and uses integer linear programming bounds to
replace the resulting affine forms.  Constant-depth integer arithmetic
then evaluates them.  It does not assume that arbitrary deep real weights
can be rounded within the original architecture.  Its depth
is $O(h+1)$, its size including wires is polynomial in the bracketed
quantity, and a linear readout of its output bits, with absolute
coefficient sum less than four, approximates the clipped target to
$\epsilon/2$.  For a threshold gate with $r$ inputs and threshold $j$,
the signed-bit affine input is
\[
 \ell(\sigma)=\sum_{i=1}^r\sigma_i-(2j-r-1).
\]
It has margin at least one on exact signs.  If the circuit size is $S$,
choose $\delta=\min\{\epsilon/16,1/[8(S+1)]\}$ and
$k=\log(2/\delta)$.  Replacing each gate by $\tanh(k\ell)$ preserves
signed-bit error at most $\delta$ at every layer: incoming error changes
$\ell$ by at most $1/8$, and $2e^{-7k/4}\le\delta$.
Copying signed bits with $\tanh(k\sigma)$ restores the same accuracy, so
arbitrary further depth causes no accumulation.

Before the tanh replacement, round the final clipped rational value to
$t=\lceil\log_2(128/\epsilon)\rceil$ binary places. If its signed
numerator is $A$ and positive denominator is $Q$, compute
$\lfloor2^t\min\{|A|,Q\}/Q\rfloor$ by the constant-depth threshold
circuit for integer division~\cite{hab}, then mask its magnitude bits
by the sign of $A$ (zero has sign zero).  This changes the value by at
most $2^{-t}$ and leaves only $2(t+1)$ potentially nonzero readout bits.
The input word lengths affect the circuit size, which is still polynomial
in the displayed resource, but not this final number of bits.  Enlarging
$b$ absorbs the additional constant-depth block.  For arbitrary additional
learner depth, copy only these final signed bits.  Their readout variation
plus bias remains less than four.

An exact constant $c_0=\tanh(1)$ can be carried by the weight $1/c_0$;
alternatively the model's affine biases supply offsets directly.
Converting the Boolean readout to signed bits gives a nonconstant
coefficient sum less than two, so its raw reconstruction error is at
most $2\delta$.  The triangle inequality proves the claimed raw output
error before any clipping.  The width is $S+1$, gate coefficients are
at most $k$, and bias/constant-wire coefficients at most $(S+1)k/c_0$.
Enlarging a numerical exponent $b$ gives all bounds.
\end{proof}

